% Capítulo 5 - Conclusão
\chapter{Conclusão}
\label{cap_conclusao}

Nos seis voos a 40~m, usar o movimento da aeronave para escolher quando rever o terreno trouxe um ganho pequeno na proporção de recipientes localizados e uma redução mais expressiva dos falsos alarmes. Com o mesmo número de recortes por voo, a memória por quadrantes localizou 25 recipientes a mais entre os 1.924 anotados e deixou 435 marcações a menos sem correspondência no mapa. Entre os 1.035 recipientes reconhecidos por ambas as estratégias, o tempo mediano entre a primeira aparição anotada e a primeira detecção correta no vídeo foi de 1,33~s com os quadrantes e de 2,37~s com o intervalo regular.

A segunda vista ajudou quando o detector deixou passar um recipiente na entrada da região observada no campo de visão. Em outras passagens, o reconhecimento ocorreu apenas na entrada, o que justifica conservar as duas oportunidades. Reunir esses alertas no mapa depende da projeção das imagens no terreno. Mesmo as caixas anotadas à mão para um mesmo recipiente apresentaram deslocamento entre as duas vistas. Nas 216 trilhas com maior variação de movimento, o intervalo regular confirmou cinco recipientes a mais que os quadrantes nos primeiros cinco segundos. Como esses mesmos seis voos orientaram a escolha da estratégia, a vantagem observada ainda precisa ser verificada em rotas independentes.

Selecionar recortes tornou a demanda compatível com a vazão medida na Jetson Orin Nano. Em bancada, com pesos treinados no COCO e imagens já decodificadas, a placa processou 44,87~recortes/s em FP16. Na simulação da fila com a vazão de bancada, todos os recortes selecionados por cada estratégia foram processados antes do fim de cada vídeo. Examinar todos os recortes permitiu localizar 321 recipientes a mais que a memória, mas exigiria cerca de 719~recortes/s para acompanhar a câmera. A simulação deixa em aberto o tempo que a decodificação do vídeo em 4K, a seleção dos recortes e a transmissão acrescentariam entre a captura e a chegada do alerta à equipe.

A altura introduziu outro compromisso. Mantidas a trajetória e a energia de cada percurso com o kit ativo, a faixa imageada seria mais larga a 80~m e a mediana estimada cairia de 270,3 para 172,3~Wh/km$^2$. Ao simular a perda de detalhe nos mesmos recortes de 40~m, a proporção de recipientes localizados pela memória caiu de 56,08\% para 46,15\%. O intervalo regular chegou a 45,95\% na escala de 80~m, praticamente o mesmo resultado. A estimativa energética inclui o terreno adicional da faixa ampliada, enquanto a avaliação do detector acompanha apenas os recipientes anotados na faixa original. Determinar quantos recipientes uma inspeção encontraria por unidade de energia a essa altura exige observar também a faixa adicional.

A estimativa do custo de levar o detector a bordo parte da energia medida pela bateria em voos sem o kit. Nos dez voos de comparação a 4,1~m/s, o erro absoluto médio do modelo de potência de cruzeiro foi de 4,2~W. A 40~m, o principal acréscimo estimado veio do transporte dos 174~g do kit, cerca de 24,8~Wh/km$^2$ na mediana. Executar o detector em FP16 acrescentaria apenas 0,33~Wh/km$^2$ ao consumo da placa já ligada. Nos cenários simulados a 40~m, reduzir a massa do kit de 250 para 100~g teria efeito maior sobre a energia por área do que trocar FP16 por INT8 no kit de 174~g.

Um ensaio em rotas independentes permitiria verificar se a antecipação dos alertas se mantém com a origem da malha fixada antes da decolagem e com a telemetria disponível no instante de cada quadro. Com o detector treinado na MBG-V2 embarcado, seria possível medir o tempo da captura até a recepção do alerta. Repetir percursos equivalentes com e sem o kit e em mais de uma altura permitiria confrontar as estimativas de energia por área com o consumo da bateria e com os recipientes efetivamente localizados.
